% Tipo de documento. En este caso es un art�culo, para folios A4, tama�o de la fuente 11pt y con p�gina separada para el t�tulo
\documentclass[a4paper,11pt,titlepage]{article}

% Carga de paquetes necesarios. OrdenesArticle es un paquete personalizado
\usepackage[spanish]{babel} 
\usepackage{eurosym} %para poder poner el s�mbolo del euro con \euro
\RequirePackage[T1]{fontenc}
\RequirePackage[ansinew]{inputenx} 
\usepackage[spanish,cap,cont,title,fancy]{OrdenesArticle}
\usepackage{lmodern}
\usepackage{graphicx}
\usepackage{hyperref}
\usepackage[usenames,dvipsnames]{color}
\usepackage{color}
\usepackage{ifthen}
\usepackage{longtable}
\usepackage{listings}
\usepackage{enumerate}
\hypersetup{bookmarksopen,bookmarksopenlevel=3,linktocpage,colorlinks,urlcolor=blue,citecolor=blue,
						linkcolor=blue,filecolor=blue,pdfnewwindow,
						pdftitle={JXTA - Aplicaciones Peer to Peer basadas en Java},
						pdfauthor={Juan Andrada Romero, Jose Domingo L�pez L�pez},
						pdfsubject={Sistemas para la Colaboraci�n}}


% Macro para definir una lista personalizada 
\newenvironment{milista}%
{\begin{list}{\textbullet}%
{\settowidth{\labelwidth}{\textbullet} \setlength{\leftmargin}{\dimexpr\labelsep+\labelwidth+5pt}
\setlength{\itemsep}{\dimexpr 0.5ex plus 0.25ex minus 0.25ex}
\setlength{\parsep}{\itemsep}
\setlength{\partopsep}{\itemsep}
\addtolength{\topsep}{-7.5pt}
}}%
{\end{list}}

% Macro para insertar una imagen
%       Uso: \imagen{nombreFichero}{Factor escala}{Caption (leyenda)}{Label (identificador para referenciarla)}
% -------------------------------------------------------------------------------------------------------------
\def\imagen#1#2#3#4{
 \begin{figure}[h]
 \begin{center}
   \scalebox{#2}{\includegraphics{#1}}
 \caption {#3}
 \label{#4}
 \end{center}
 \end{figure}
}

% Definicion del "listings" para el lenguaje Java
\lstdefinelanguage{Java}
{
 morecomment = [l]{//}, 
 morecomment = [l]{///},
 morecomment = [s]{/*}{*/},
 morestring = [b]", 
 sensitive = true,
 morekeywords = {package, static, while, switch, break, line, void, String, Object, int, Integer, instanceof, else, if, for, private, return, new, public, class, import, int, boolean, true, false, extends, final, super, protected, abstract, this, do, float, double, null, try, catch, implements}
}

% Configuraci�n del c�digo incluido para el lenguaje Java
\lstset{
  language=Java,
  basicstyle=\footnotesize,
  backgroundcolor=\color{white},
  showspaces=false,
  showstringspaces=false,
  showtabs=false,
  frame=single,
  tabsize=2,
  captionpos=b,
  breaklines=true,
  breakatwhitespace=false,
  escapeinside={\%},
  keywordstyle = \color [rgb]{0,0,1},
  commentstyle = \color [rgb]{0.133,0.545,0.133},
  stringstyle = \color [rgb]{0.627,0.126,0.941}
}


\begin{document}

% En las p�ginas de portada e �ndices, no hay encabezado ni pie de p�gina
\pagestyle{empty} 

% Se incluye la portada
\begin{titlepage}
	\begin{center}
  	{\LARGE UNIVERSIDAD DE CASTILLA-LA MANCHA} \\
  	\bigskip
  	{\Large ESCUELA SUPERIOR DE INFORM�TICA} \\
  	\vspace{20mm}
  	\includegraphics[scale=0.45, keepaspectratio]{./images/esi_bw.png} \\
  	\vspace{20mm}
  	{\Huge \textbf{Almacenamiento y Recuperaci�n de la Informaci�n}}\\
  	\vspace{10mm}
  	{\LARGE \textsc{\textbf{M�dulo de importaci�n e indexaci�n de documentos}}}\\
  	\vspace{30mm}
  	\today
  	\vspace{40mm}
  	\begin{flushleft}
  		{\large Juan Andrada Romero}\\
  		\vspace{1mm}
  		{\large Jose Domingo L�pez L�pez}\\		
  	\end{flushleft}
	\end{center}
\end{titlepage}
\clearpage

% Texto del reverso de la portada
%%\mbox{}
%%\vspace{18cm}
%%\begin{small}
% Se ajusta la separaci�n entre p�rrafos
%%\parskip=10pt 

%%\copyright~ 2008/2009 Juan Andrada Romero. Universidad de Castilla La Mancha, Escuela Superior de Inform�tica de Ciudad Real.

%%Se permite la modificaci�n, copia y distribuci�n de este documento, seg�n la licencia de documentaci�n GNU (\url{http://www.gnu.org}).

%%Este documento fue compuesto con \LaTeX{}. Im�genes generadas con Power Point y Gimp.
%%\end{small}

%%\newpage

% En las p�ginas de �ndices y prefacio, se utiliza numeraci�n romana
\pagenumbering{Roman}

% Se crea el �ndice
\tableofcontents
% Se pasa p�gina y se a�ade el �ndice de figuras al �ndice principal
%\clearpage\phantomsection
%\addcontentsline{toc}{section}{\listfigurename}
% Se crea el �ndice de figuras
%\listoffigures
% Si se quiere crear un �ndice de tablas se pondr�a: \listoftables

\newpage

% Se ajusta la separaci�n entre p�rrafos
\parskip=10pt

% Se a�ade el prefacio al �ndice
%\clearpage\phantomsection
%\addcontentsline{toc}{section}{Prefacio}

% Comienza el contenido del documento. Se utilizan n�meros ar�bigos y el encabezado y pie de p�gina personalizado
\pagenumbering{arabic}
\pagestyle{fancy}

\section{Introducci�n}


\clearpage
\section{Conceptos de JXTA}


\subsection{Par (\textit{Peer})}

Un par (\textit{peer)} es cualquier entidad de red que implementa uno o m�s protocolos utilizados en JXTA. As�, un par puede residir en una PDA, un tel�fono m�vil, un PC, etc, trabajando de manera independiente y as�ncrona con respecto a otros pares.

Un par debe publicarse en la red para que pueda ser descubierto por otros pares y puedan establecer una comunicaci�n. Por tanto, cada par va a tener un identificador �nico y una direcci�n de red, conocida esta �ltima como \textit{peer endpoint} y utilizada para establecer una comunicaci�n punto a punto entre dos pares. \\
\indent Normalmente, los pares se configuran para ser descubiertos espont�neamente por otros pares de la red, a partir del concepto de \textit{grupo de pares}, que puede ser, por naturaleza, persistente o transitorio.

Los pares en JXTA pueden clasificarse en tres categor�as:

\begin{milista}
	\item \textbf{\textit{Minimal-Edge Peers}}: pares que implementan s�lo los servicios principales JXTA, utilizando otros pares para que act�en como su proxy a la hora de utilizar otros servicios que no pertenecen al \textit{core}de JXTA. Pares que pertenezcan a esta categor�a pueden ser sensores.
	\item \textbf{\textit{Full-Edge Peers}}: pares que implementan todos los servicios de JXTA y pueden participar en todos sus protocolos. Estos pares conforman la mayor�a de pares en una red JXTA y pueden ser tel�fonos, PC's, servidores, etc.
	\item \textbf{\textit{Super Peer}}: pares que implementan y proveen recursos para soportar el despliegue y operaci�n de una red JXTA. Hay tres funciones principales que debe tener un \textit{Super-Peer}:
	\begin{itemize}
		\item \textbf{Retransmisi�n}: usada para almacenar y reenviar mensajes entre pares que no tienen conectividad directa debido a cortafuegos o redes privadas.
		\item \textbf{Citas}: se ocupa de mantener los anuncios globales sobre donde est�n publicados otros pares y ayuda a los pares que act�an como proxy para buscar esos anuncios. 
		\item \textbf{Proxy}: utilizado por los \textit{minimal-edge peers} para tener acceso a todas las funcionalidades de la red JXTA. El par proxy traduce las peticiones, responde a las consultas y proporciona soporte para la funcionalidad de los \textit{minimal-edge peers}.
	\end{itemize}

\end{milista}


\subsection{Grupo de pares (\textit{Peer Group})}

Un grupo de pares es una colecci�n de pares que comparten un conjunto com�n de servicios o intereses. Los pares se auto-organizan en grupos, donde cada uno de los grupos se identifica por un identificador �nico. Cada grupo establece su propia pol�tica de pertenencia al grupo, que puede variar desde una pol�tica abierta (cualquier par puede asociarse al grupo) hasta la muy segura y protegida (que requiere credenciales para poder ser miembro).

Los pares pueden pertenecer a m�s de un grupo simult�neamente. Por defecto, el primer grupo que se instancia es el \textit{Network Peer Group}, y, posteriormente, todos los pares pueden optar por unirse a otros grupos en cualquier momento.

Los protocolos JXTA describen como los pares pueden publicarse, descubrirse, unirse, y monitorear grupos de pares; sin embargo, dichos protocolos no indican cu�ndo o por qu� se crean los grupos de pares. Por otra parte, unirse a un grupo consiste simplemente en instanciar todos los servicios definidos por ese grupo. \\
\indent Hay varias motivaciones para la creaci�n de grupos de pares:

\begin{milista}
	\item Para crear un entorno seguro, ya que se puede aplicar una pol�tica de seguridad. �sta puede ser tan simple como un intercambio de texto plano con nombre de usuario y contrase�a, o tan sofisticada como criptograf�a de llave publica. As�, los miembros del grupo pueden publicar y acceder a contenido protegido.
	\item Para delimitar el alcance, ya que los grupos que permiten la creaci�n de un dominio local espec�fico. Por ejemplo, los pares pueden agruparse para implementar una red de documentos compartidos o una red de CPUs compartidas. Los grupos de pares sirven para subdividir la red en regiones abstractas proporcionando un mecanismo impl�cito de alcance. De este modo, las fronteras del grupo definen el alcance de la b�squeda cuando se busca algo dentro del grupo.
	\item Para crear un entorno de monitorizaci�n, ya que permiten monitorear un conjunto de pares para un fin especial.
\end{milista}


\subsection{Servicios de red (\textit{Network Services})}

Los pares cooperan y se comunican para publicar, descubrir, e invocar servicios de red. Adem�s, los pares pueden publicar m�ltiples servicios los cu�les, a su vez, son descubiertos a trav�s del \textit{Peer Discovery Protocol}. De este modo, los protocolos JXTA reconocen dos niveles de servicios de red:

\begin{milista}
	\item \textbf{Servicios de par}: un servicio de este tipo es solamente accesible por el par que est� publicando ese servicio. Si ese par fallara, el servicio tambi�n fallar�a. 
	\item \textbf{Servicios de grupo de pares}: un servicio de grupo de pares est� compuesto de una colecci�n de instancias del servicio, cooperando entre s� y funcionando en m�ltiples miembros del grupo. S� falla cualquiera de los pares, el servicio colectivo del grupo no se ver� afectado siempre y cuando el servicio este a�n disponible desde al menos un par miembro. 
\end{milista}


\subsection{Servicios de grupos de pares (\textit{Peer Group Services})}

JXTA define un conjunto b�sico de servicios de grupos de pares, que todos los pares miembros del grupo deben implementar, al igual que los pares que se unan a dicho grupo. As�, para que dos pares puedan interactuar a trav�s de un servicio, deben ser ambos miembros del mismo grupo. Los servicios de grupo que cada uno de los pares debe implementar con el fin de participar en una red JXTA son:
 
\begin{milista}
	\item\textbf{Servicio de Destino (\textit{Endpoint Service})}: el servicio de destino o de \textit{punto final} es utilizado para enviar y recibir mensajes entre pares. Este servicio normalmente implementa el protocolo de encaminamiento al destino (ERP), aunque implementaciones m�nimas proporcionan s�lo una peque�a parte del total de este protocolo.
	\item\textbf{Servicio de Resoluci�n (\textit{Resolver Service})}: el servicio de resoluci�n es utilizado para enviar solicitudes de consulta gen�rica a otros pares. Los pares pueden definir e intercambiar consultas para encontrar cualquier informaci�n que pueda ser necesaria (por ejemplo, encontrar otros pares, determinar el estado de un servicio o la disponibilidad de una tuber�a, etc.).
\end{milista}

Adem�s de los servicios b�sicos anteriores, se pueden definir otros servicios est�ndares en los grupos, aunque �stos no tienen por qu� ser implementados por todos los grupos, sino s�lo por aquellos que los necesiten. Los servicios est�ndares que se encuentran en la mayor�a de grupos de pares son:  
 
\begin{milista}

	\item\textbf{Servicio de Descubrimiento (\textit{Discovery Service})}: el servicio de descubrimiento es usado por los pares miembros para buscar recursos en el grupo, tales como otros pares, otros grupos, servicios y tuber�as. 
	\item\textbf{Servicio de afiliaci�n a un grupo(\textit{Membership Service})}: es utilizado por los pares para identificarse con seguridad dentro de un grupo. La identidad de un par permite a las aplicaciones y los servicios determinar qui�n est� solicitando una operaci�n y decidir si la operaci�n debe permitirse o no. Las aplicaciones pueden realizar su propio control de acceso o pueden utilizar el servicio de acceso.
	\item\textbf{Servicio de Acceso (\textit{Access Service})}: el servicio de acceso es usado para validar las peticiones hechas por un par a otro. El par que recibe la petici�n verifica las credenciales del par que la solicita para determinar si se le permite el acceso. Sin embargo, no todas las acciones dentro del grupo necesitan comprobarse con el servicio de acceso; s�lo aquellas acciones que son limitadas a algunos pares necesitan ser comprobadas.
	\item\textbf{Servicio de Tuber�as (\textit{Pipe Service})}: el servicio de tuber�a es utilizado para crear y manejar conexiones de tuber�a  entre los miembros del grupo.
	\item\textbf{Servicio de Monitoreo (\textit{Monitoring Service})}: el servicio de monitoreo se utiliza para permitir a un par monitorizar a otros miembros del mismo grupo de pares. 

\end{milista}


\subsection{M�dulos}

Los m�dulos de JXTA son una abstracci�n de JXTA de bajo nivel usada para representar cualquier pieza de ''c�digo'' y el API que proporciona ese c�digo para su uso. Los m�dulos son utilizados para implementar en JXTA servicios, transporte de mensajes, etc.
 
Los m�dulos proveen una abstracci�n gen�rica para permitir a un par instanciar una funci�n o servicio como, por ejemplo, cuando se une a un grupo, ya que necesitan implementar los servicios de ese grupo. Los m�dulos permiten la representaci�n y la publicaci�n de funcionalidades independientes de la plataforma, y permite a los pares describir e instanciar cualquier tipo de implementaci�n de una funcionalidad. Por ejemplo, un par puede instanciar una implementaci�n en Java o C de una funcionalidad espec�fica.
 
La posibilidad para describir y publicar (o anunciar) una funcionalidad independiente de la plataforma es esencial para soportar el desarrollo de nuevos servicios de grupo. De hecho, JXTA utiliza los anuncios de m�dulos para auto-describir sus propios servicios.

Un m�dulo incluye los siguientes elementos: 

\begin{milista}
	
	\item \textbf{Clases}: las clases de un m�dulo se utilizan principalmente para anunciar la existencia de una funcionalidad. Cada clase de un m�dulo es identificada por un �nico ID, el \textit{ModuleClassID}.
	
	\item \textbf{Especificaci�n}: la especificaci�n de un m�dulo se utiliza para acceder a un m�dulo, ya que contiene toda la informaci�n necesaria para acceder o invocar dicho m�dulo. Por tanto, la especificaci�n de un m�dulo es un m�todo para proporcionar la funcionalidad que implica ese m�dulo, pudiendo existir diversas especificaciones para un determinada clase. Cada especificaci�n se identifica de manera �nica, a trav�s del \textit{ModuleSpecID}.

	\item \textbf{Implementaci�n}: la implementaci�n de un m�dulo es la implementaci�n de una determinada especificaci�n, pudiendo existir m�ltiples implementaciones para una determinada especificaci�n del m�dulo. Cada implementaci�n del m�dulo contiene el \textit{ModuleSpecID} de la especificaci�n asociada que implementa.

\end{milista}

As�, los servicios JXTA pueden usar los m�dulos para identificar la existencia del servicio (las clases del m�dulo), la especificaci�n del servicio (una especificaci�n del m�dulo), o una implementaci�n del servicio (una implementaci�n del m�dulo). Cada uno de estos componentes tiene un anuncio asociado el cu�l puede ser publicado y descubierto por otro par JXTA.
 
Por ejemplo, el servicio de descubrimiento de JXTA (\textit{Discovery Service}) tiene un �nico \textit{ModuleClassID}, identific�ndolo como un servicio de descubrimiento, que es su funcionalidad abstracta. Adem�s, puede tener m�ltiples especificaciones del servicio, utilizando diferentes estrategias para realizar el descubrimiento, adaptadas al tama�o del grupo y su dispersi�n a trav�s de la red. Para cada especificaci�n, puede haber m�ltiples implementaciones, cada una de las cuales contiene el mismo \textit{ModuleSpecID}.


\subsection{Mensajes}

Los servicios y aplicaciones JXTA se comunican usando mensajes JXTA, que son las unidades b�sicas de intercambio de datos entre pares. Cada protocolo JXTA est� definido como un conjunto de mensajes que intercambian los pares participantes. Los mensajes son enviados entre los pares utilizando el servicio de destino (\textit{Endpoint Service}), el servicio de tuber�a (\textit{Pipe Service}), as� como un socket y otras v�as de acceso. La mayor�a de las aplicaciones no necesitan utilizar directamente tuber�as unidireccionales o el servicio de destino, sino que utilizan com�nmente \textit{sockets} de comunicaci�n y tuber�as bidimensionales.
 
Los protocolos de JXTA est�n especificados como un conjunto de mensajes intercambiados entre pares. El uso de mensajes de XML para definir los protocolos permite a muchos tipos diferentes de pares utilizar un protocolo determinado. Debido a que los datos est�n etiquetados, cada par es libre para implementar el protocolo de un modo m�s adecuado a sus capacidades y rol. Por ejemplo, un par que est� muy limitado y no tenga capacidad suficiente para procesar la totalidad de un mensaje, puede usar las etiquetas de datos para extraer las partes que pueda procesar e ignorar las restantes. 

\subsection{Tuber�as (\textit{Pipes})}

Los pares de JXTA utilizan tuber�as para enviar mensajes a otro par. Por tanto, las tuber�as son un mecanismo de transferencia de mensajes as�ncrono, unidireccional y no confiable (con la excepci�n de las tuber�as punto a punto segurass) usados para la comunicaci�n y transferencia de datos. Las tuber�as son canales de comunicaci�n virtuales y pueden conectar pares que no tienen un vinculo f�sico directo, lo que se traduce en una conexi�n l�gica. Las tuber�as pueden ser usadas para enviar cualquier tipo de datos incluyendo texto XML y HTML, im�genes, m�sica, c�digo binario, cadenas de datos y objetos Java.

La tuber�as de destino o de ''punto final'' (\textit{endpoints}) est�n din�micamente vinculadas con los pares de destino (correspondientes a interfaces de red, como un puerto TCP) cuando la tuber�a se abre. Toda la resoluci�n y comunicaci�n de tuber�as se realiza dentro del alcance de un grupo de pares. es decir, las tuber�as de salida y entrada deber�n pertenecer al mismo grupo de pares.
 
Las tuber�as ofrecen dos modos de comunicaci�n: punto a punto (\textit{unicast}) y propagaci�n, como se puede observar en la Figura \ref{fig:pipes}.

\imagen{./imagenes/tuberias}{0.7}{Tipos de tuber�as}{fig:pipes}
 
\begin{milista}
	\item \textbf{Tuber�as punto a punto}: una tuber�a punto a punto conecta exactamente dos tuber�as de punto final. La tuber�a de entrada en un par recibe mensajes enviados desde la tuber�a de salida del otro par. Tambi�n es posible para varios pares unirse a una sola tuber�a de entrada.
 	\item \textbf{Tuber�as de propagaci�n}: este tipo de tuber�as conecta una tuber�a de salida a varias tuber�as de entrada. Los mensajes fluyen desde la tuber�a de salida, que es la fuente de propagaci�n, hacia las tuber�as de entrada.
 	\item \textbf{Tuber�as punto a punto seguras}: es un tipo de tuber�a punto a punto que provee un canal de comunicaci�n seguro y fiable.
\end{milista}
 

\subsection{Anuncios (\textit{Advertisements})}

Todos los recursos de la red JXTA (servicios, pares, grupos de pares y tuber�as) son representados como anuncios. Los anuncios son un lenguaje de estructuras de meta-datos representados como documentos XML (ver Figura \ref{fig:anuncio}). Los protocolos JXTA usan anuncios para describir y publicar la existencia de recursos de un par. Del mismo modo, los pares descubren recursos mediante la b�squeda de sus correspondientes anuncios y pueden guardar en cach� cualquier anuncio local descubierto.
 
Cada anuncio es publicado con un tiempo de vida que especifica la disponibilidad de sus recursos asociados. El tiempo de vida permite la supresi�n de los recursos obsoletos sin la necesidad de un control centralizado, de modo que un anuncio puede ser publicado de nuevo (antes de que el anuncio original expire) para extender la vida �til de un recurso.
 
Los protocolos JXTA definen los siguientes tipos de anuncios:

\begin{milista}
	\item \textbf{Anuncios de pares (\textit{Peer Group Advertisement})}: describen los recursos de un par. El principal uso de estos anuncios es para mantener informaci�n espec�fica sobre el par, como su nombre, su identificador, destinos disponibles, y cualquier atributo que se desee publicar en tiempo de ejecuci�n.
	\item \textbf{Anuncios de grupos de pares (\textit{Peer Group Advertisement})}: describe los recursos de un grupo espec�fico de pares, como son el nombre, el identificador del grupo, descripci�n, etc.
	\item \textbf{Anuncios de tuber�a (\textit{Pipe Advertisement})}: describe un canal de comunicaci�n de una tuber�a, y es usado por el servicio de tuber�a para crear las tuber�as asociadas a una entrada y a una salida. Cada anuncio de tuber�a contiene un identificador simb�lico y un tipo de tuber�a (punto a punto, propagada, segura). 
		\item \textbf{Anuncio de clases de un m�dulo (\textit{Module Class Advertisement})}: describe una clase de un m�dulo. Su objetivo principal es documentar formalmente la existencia de una clase del m�dulo. Incluye un nombre, descripci�n, y un �nico identificador (\textit{ModuleClassID}).
		\item \textbf{Anuncio de especificaci�n de un m�dulo (\textit{Module Specification Advertisement})}: describe una especificaci�n de un m�dulo. Su principal objetivo es proveer referencias a la documentaci�n necesaria con el fin de crear implementaciones conforme con la especificaci�n. Un uso secundario es hacer correr instancias que se puedan utilizar de manera remota mediante la publicaci�n de informaci�n. Este anuncio incluye nombre, descripci�n, identificador �nico de la especificaci�n del m�dulo (\textit{ModuleSpecID}), anuncio de tuber�a y un campo que contiene par�metros arbitrarios a ser interpretados por cada aplicaci�n.
	\item \textbf{Anuncios de implementaci�n de un m�dulo (\textit{Module Implementation Advertisement})}: define una implementaci�n de una especificaci�n de un m�dulo. Incluye un nombre, el \textit{ModuleSpecID} asociado, as� como el c�digo, paquete y campos de par�metros que permiten a un par recuperar los datos necesarios para ejecutar la implementaci�n.
	\item \textbf{Anuncio de citas (\textit{Rendezvous Advertisement})}: describe como se establecen citas entre pares de un grupo.
	\item \textbf{Anuncio de informaci�n de un par (\textit{Peer Info Advertisement})}: describe los recursos de informaci�n de un par. El uso principal de este anuncio es contener informaci�n especifica acerca del estado actual de un par, tales como el tiempo de actividad, cuenta de mensajes de entrada y salida, el tiempo que ha pasado desde el ultimo mensaje recibido y el tiempo que ha pasado desde el ultimo mensaje enviado.
\end{milista}
  
\imagen{./imagenes/anuncio}{0.6}{Ejemplo de un anuncio de una tuber�a}{fig:anuncio}
 

\subsection{Seguridad}

Redes P2P din�micas, como la red JXTA, necesitan soportar diferentes niveles de acceso a los recursos. En JXTA, los pares operan en un modelo basado en roles, en el cual un par individual act�a bajo la autoridad de otro par de confianza para realizar una tarea.

Se proveen cinco requisitos b�sicos de seguridad en JXTA: 

\begin{milista}
	\item \textbf{Confidencialidad}: garantiza que el contenido de un mensaje no se divulgar� a pares no autorizados.
	\item \textbf{Autenticaci�n}: garantiza que el emisor de un mensaje es qui�n dice ser.
	\item \textbf{Autorizaci�n}: garantiza que el emisor est� autorizado para enviar un mensaje.
	\item \textbf{Integridad de los datos}: garantiza que el mensaje no ha sido modificado accidental o deliberadamente mientras se transmit�a.
	\item \textbf{Refutabilidad}: garantiza que el mensaje fue transmitido por un emisor identificado y que no es una respuesta de un mensaje previamente transmitido.
\end{milista}

Los mensajes XML permiten a�adir nuevos meta-datos como credenciales, certificados y claves p�blicas a los mensajes de JXTA, permitiendo que se cumplan los requisitos de seguridad anteriores. Adem�s, los mensaje se pueden encriptar y firmar.

Un credencial es un token que se utiliza para identificar un emisor y para comprobar que tiene autorizaci�n para enviar el mensaje. Adem�s, este token est� presenta en cada mensaje que ese emisor env�a, comprob�ndose con la identidad del emisor cuando se recibe el mensaje. Cabe destacar que la implementaci�n de las credenciales se realiza como una configuraci�n de un plug-in, por lo que pueden exisitr m�ltiples configuraciones de credenciales dentro de la misma red. Con esto se intenta que JXTA sea compatible con otros mecanismos de seguridad como SSL (\textit{Secure Socket Layer)} e IPSec (\textit{Internet Protocol Secure}).
\clearpage
\section{Arquitectura Software de JXTA}

En los siguientes apartados se va a detallar la arquitectura software de JXTA y los componentes que en ella intervienen.

\subsection{Descripci�n}

La arquitectura software de JXTA est� dividida en tres capas, como se muestra en la Figura \ref{fig:arquitecturaJXTA}.

\imagen{./imagenes/arquitecturaJXTA}{0.21}{Arquitectura de JXTA}{fig:arquitecturaJXTA}

\begin{milista}
	\item \textbf{N�cleo (\textit{core}) de JXTA}: el n�cleo de JXTA encapsula las primitivas m�nimas y esenciales que se utilizan com�nmente para la comunicaci�n P2P. Esto incluye componentes b�sicos para permitir las funcionalidades clave para las aplicaciones P2P, incluyendo descubrimiento, la creaci�n de pares y grupos de pares y primitivas asociadas con seguridad.
 
 \item \textbf{Capa de servicios}: la capa de servicios incluye los servicios de red que podr�an no ser absolutamente necesarios para que una red P2P opere, pero que son comunes o deseables. Los ejemplos de servicios de red incluyen b�squeda e indexaci�n, directorios, sistemas de almacenamiento, uso compartido de archivos, sistemas de archivos distribuidos, agregaci�n y petici�n de recursos, autenticaci�n y servicios de PKI (\textit{Public Key Infrastructure}).
 
\item \textbf{Capa de aplicaciones}: la capa de aplicaciones incluye la implementaci�n de aplicaciones integradas, como aplicaciones P2P de  mensajer�a instant�nea, distribuci�n de documentos y de recursos, sistemas de correo electr�nico P2P, sistemas de subasta distribuido, etc.

\end{milista}

Sin embargo, la frontera entre aplicaciones y servicios no es r�gida, ya que una aplicaci�n de un cliente puede considerarse como un servicio para otro cliente. El sistema JXTA est� dise�ado para ser modular, permitiendo a los desarrolladores seleccionar y escoger una colecci�n de servicios y aplicaciones que se adapten a sus necesidades.


\subsection{Componentes de JXTA}

Una red JXTA consiste en una serie de nodos (pares) interconectados mediante cualquier protocolo de interconexi�n adecuado, incluyendo TCP/IP, HTTP, Bluetooth, GSM, etc. Adem�s, m�ltiples pares pueden correr sobre un �nico dispositivo f�sico y varios dispositivos f�sicos podr�an cooperar para actuar como un solo par.

Cada par ofrece un conjunto de servicios y recursos que est�n a disposici�n de otros pares. As�, existen dos tipos de servicios comunes en las redes JXTA, que son los servicios de pares y servicios de grupo. De este modo, todos los pares JXTA implementan un peque�o n�mero de servicios b�sicos necesarios y tambi�n proporcionan varios servicios est�ndares adicionales. Cada grupo de pares, por su parte, incluye el conjunto de servicios de grupo que cada par debe ejecutar para participar en ese grupo.

Los recursos de un par son normalmente contenido est�tico (no-interactivo) los cu�les �l controla, posee o simplemente tiene una copia. Los recursos pueden incluir archivos, documentos, �ndices, pero tambi�n pueden ser recursos del mundo real tales como switches, sensores e impresoras. Dichos servicios y recursos son publicados como un \textit{anuncio}.

Otro componente necesario en una red JXTA son \textit{sockets} y tuber�as, que se utilizan para enviar mensajes entre pares. Como ya se ha comentado en la secci�n \ref{sec:conceptos}, los \textit{sockets} de JXTA son conexiones bidireccionales confiables usadas por las aplicaciones para comunicarse., y las tuber�as son un mecanismo de transferencia de mensajes as�ncrono y unidireccional usado para la comunicaci�n del servicio. Los mensajes son simples documentos XML que contienen informaci�n acerca de la ruta a seguir, un resumen y credenciales. 

Para terminar, se destacan cuatro aspectos esenciales de la arquitectura JXTA que la distinguen de otros modelos de red distribuida.

\begin{enumerate}
	\item El uso de documentos XML (\textit{anuncios}) para describir recursos de red.
	\item La abstracci�n de tuber�as a pares y de pares a ''puntos finales'' (\textit{endpoints}), sin depender de una autoridad central de nombramiento/direccionamiento, tal como DNS.
	\item Un esquema de direccionamiento uniforme de pares.
	\item Una infraestructura de b�squeda descentralizada basada sobre tablas Hash distribuidas (\textit{Distributed Hash Table - DHT}) para la  indexaci�n de recursos.
\end{enumerate}


\clearpage
\section{Arquitectura de red de JXTA}

\subsection{Organizaci�n de la red}

La red JXTA es una red \textit{Ad-Hoc}, multi-salto y adaptativa, compuesta de pares conectados. Las conexiones en la red pueden ser temporales y, como resultado, el enrutamiento de mensajes entre pares es no determinista, ya que los pares pueden unirse o dejar la red en cualquier momento, lo que resulta siempre en un cambio de informaci�n de enrutamiento.

Normalmente se utilizan cuatro tipos de pares en una red JXTA, como se observa en la Figura \ref{fig:redJXTA}:

\begin{milista}
	\item \textbf{\textit{Minimal-Edge Peer}}: puede enviar y recibir mensajes, pero no guardar anuncios en cach� o enrutar mensajes para otros pares. 
	\item \textbf{\textit{Full-featured Peer}}: puede enviar y recibir mensajes y t�picamente guardara en cach� los anuncios. Este tipo de par simple contesta a peticiones de descubrimiento con informaci�n encontrada en su cach� de anuncios, pero no reenv�a ninguna petici�n de descubrimiento. 
	\item \textbf{\textit{Rendezvous peer}}: es un par de infraestructura que ayuda a otros pares con la propagaci�n de mensajes, descubrimiento de anuncios y rutas y mantiene un mapa topol�gico de otros pares de infraestructura, que se utilizan para controlar la propagaci�n y el mantenimiento de la tabla Hash distribuida. Cada grupo de pares mantiene su propio conjunto de pares \textit{rendezvous} y puede tener tantos como sea necesario. 
	\item \textbf{\textit{Relay peer}}: es un par de infraestructura que ayuda a otros pares no direccionables con la retransmisi�n del mensaje. Un par puede solicitar una secci�n del mensaje a este tipo de par para facilitar la retransmisi�n del mensaje cada vez que sea necesario.
\end{milista}

\imagen{./imagenes/arquitecturaRed}{0.4}{Arquitectura de red de JXTA}{fig:redJXTA}


\subsection{�ndice Distribuido de Recursos Compartidos (\textit{Shared Resource Distributed Index - SRDI})}

JXTA implementa un SRDI para conseguir un mecanismo eficiente de propagaci�n de peticiones en la red JXTA. Los pares de tipo \textit{rendezvous} mantienen un �ndice de anuncios publicados por los pares de tipo \textit{edge}, de tal modo que cu�ndo uno de estos pares publica un anuncio, el par \textit{rendezvous} lo indexa en su SRDI, utilizando como clave el ID del anuncio.

Cada par de tipo \textit{rendezvous} conoce a otros pares de este tipo que se encuentren en su grupo, de tal modo que cuando se indexa un nuevo recurso, el �ndice utilizado se env�a al resto de pares \textit{rendezvous}, para que tambi�n lo almacenen en su SRDI. 

A continuaci�n, se va a comentar un ejemplo de uso del SRDI cuando un par solicita un recurso, que es el mostrado en la Figura \ref{fig:ejSRDI}:

\imagen{./imagenes/ejemploSRDI}{0.2}{Propagaci�n de una petici�n en una red JXTA usando SRDI}{fig:ejSRDI}

\begin{milista}
	\item El par A de tipo \textit{edge} solicita un recurso a su par R1 de tipo \textit{rendezvous} y al resto de la red, usando difusi�n (\textit{multicast}).
	\item Si los pares de la red de A tienen el recurso en su cach�, se lo env�an.
	\item Del mismo modo, el par R1 usa el �ndice del anuncio de petici�n para buscarlo en su SRDI. Si lo encuentra, se lo env�a al par A. Si no lo encuentra, propaga la petici�n a los pares \textit{rendezvous} que se encuentran en su grupo, para que �stos realicen la b�squeda o repitan este proceso si no tienen ese recurso disponible.
\end{milista}


\subsection{Firewalls y NAT}

Un par detr�s de un cortafuegos (\textit{firewall}) puede enviar un mensaje directamente a un par que est� fuera del cortafuegos, pero no es posible establecer una comunicaci�n directa en el otro sentido. Lo mismo ocurre para los pares que se encuentren tras un dispositivo de NAT (\textit{Network Address Translation}).

Para poder realizar la conexi�n con pares que est�n tras cortafuegos o usan NAT, se deben cumplir las siguientes condiciones:

\begin{milista}
	\item Al menos un par del grupo que se encuentra en el cortafuego debe poder comunicarse con un par tras el cortafuegos.
	\item Estos pares deben soportar un protocolo com�n de transporte, como TCP o HTTP.
	\item El cortafuegos debe permitir tr�fico HTTP o TCP. 
\end{milista}


	
	
\clearpage
\section{Protocolos utilizados en JXTA}

JXTA define una serie de mensajes XML, o protocolos, para la comunicaci�n entre pares. Los pares utilizan estos protocolos para descubrir uno a otro, anunciar y descubrir los recursos de red.

Existen seis tipos de protocolos en JXTA, que son as�ncronos y se basan en un modelo de consultas y respuestas. Un par JXTA utiliza uno de los protocolos para enviar una consulta a uno o m�s pares de su grupo. Es posible que reciba cero, una o m�s respuestas a su consulta. Por ejemplo, un par puede utilizar el protocolo PDP para enviar una consulta de descubrimiento, solicitando todos los pares registrados en el grupo por defecto (\textit{Net Peer Group}). En este caso, varios pares responder�n con sus respuestas de descubrimiento. \\
\indent Los pares JXTA no necesitan implementar todos los protocolos, s�lo necesitan implementar los protocolos que vayan a utilizar. 

En los siguientes apartados se detallan cada uno de los protocolos que existen en JXTA.
 
 
\subsection{Protocolo de Descubrimiento de Pares (\textit{PDP - Peer Discovery Protocol})}

El protocolo de descubrimiento de pares (PDP) se utiliza para descubrir cualquier recurso publicado por otros pares en forma de anuncio. Un recurso puede ser un par, un grupo de pares, una tuber�a, un servicio o cualquier otro recurso que tenga un anuncio.

PDP es el protocolo de descubrimiento por defecto para todos los grupos de pares definidos por el usuario y para el el grupo de pares por defecto. De este modo, si se definen servicios de descubrimiento personalizado, �stos se pueden aprovechar del PDP, redefiniendo las partes que sean necesarias.

Existen varias formas para descubrir informaci�n distribuida. La implementaci�n actual de JXTA utiliza una combinaci�n de propagaci�n (\textit{multicast}) IP a la subred local y el uso de una red de citas, que es una t�cnica basada en citas mantenido la DHT (tabla hash distribuido). Los nodos de la red de tipo \textit{rendezvous} proporcionan el mecanismo necesario para dirigir las solicitudes y poder descubrir de forma din�mica la informaci�n. Adem�s, un nodo puede configurarse para que, al iniciarse, localice din�micamente nodos o recursos de su red local a trav�s de mensajes de multidifusi�n.

\subsection{Protocolo de Informaci�n de Pares (\textit{PIP - Peer Information Protocol})}

Este protocolo es utilizado por los pares para obtener informaci�n de estado (el tiempo de actividad, estado, reciente de tr�fico, etc.) de otros pares. As�, una vez que se encuentra un nodo utilizando el protocolo anterior, se pueden consultar sus capacidades y su estado. 

El protocolo de informaci�n de pares proporciona un conjunto de mensajes para obtener informaci�n del estado de un par. Esta informaci�n puede utilizarse para la implementaci�n comercial o interna de aplicaciones JXTA. Por ejemplo, en las implementaciones comerciales, la informaci�n puede utilizarse para determinar el uso de un servicio de pares y facturar a los consumidores por el uso de ese servicio. En una implementaci�n interna de una empresa, la informaci�n puede utilizarse por el departamento de tecnolog�as de la informaci�n (TI) para supervisar el tr�fico de red para mejorar el rendimiento general.

Entre los mensajes que permite utilizar este protocolo, destacan el mensaje de \textit{ping} y el mensaje de \textit{PeerInfo}. El primero de ellos se env�a a un para comprobar est� vivo y para obtener informaci�n acerca de �l. El mensaje de ping especifica si se devuelve una respuesta completa (anuncio de pares) o un mensaje de notificaci�n (si est� vivo o no y tiempo de actividad).\\
\indent  Por otra parte, el mensaje de \textit{PeerInfo} se utiliza para enviar un mensaje en respuesta a un mensaje de \textit{ping}. Contiene las credenciales del remitente, el ID de los pares de origen y destino, el tiempo de actividad y el anuncio de pares. 

\subsection{Protocolo de Resoluci�n de Pares (\textit{PRP - Peer Resolver Protocol})}

El protocolo de resoluci�n de pares (PRP) permite a un par enviar solicitudes de consultas gen�ricas a otros pares, e identificar respuestas coincidente a esas consultas. Las solicitudes de consulta pueden ser enviadas a un punto espec�fico o pueden ser propagadas a trav�s de los servicios de cita a otros pares en el �mbito de un grupo de pares. De este modo, el PRP utiliza el servicio de citas para difundir una consulta a varios pares y utiliza mensajes \textit{unicast} para enviar consultas a pares espec�ficos.

El PRP es un protocolo base para apoyar las solicitudes de consultas gen�ricas. Tanto el PIP como el PDP se construyen usando PRP y proporcionan consultas y peticiones espec�ficas: el PIP se utiliza para consultar la informaci�n de estado y el PDP se utiliza para descubrir recursos de pares. Para cualquier consulta gen�rica que puede ser necesaria en una aplicaci�n, se puede utilizar el PRP. 

El mensaje de respuesta a la consulta (mensaje de resoluci�n de la consulta) contiene la credencial del par que env�a la petici�n, un controlador de servicio espec�fico, un ID �nico de la consulta y el remitente. Cada servicio puede registrar un controlador en el grupo de pares para procesar las solicitudes de consulta y generar las respuestas.


\subsection{Protocolo de Enlace a Tuber�as (\textit{PBP - Pipe Binding Protocol})}

El protocolo de enlace de tuber�a (es utilizado por miembros del grupo de pares para enlazar un anuncio de tuber�a a un par que se encuentra en el extremo de la tuber�a. De este modo, una tuber�a puede ser vista como una cola de mensajes con nombre abstracto, que soporta operaciones de creaci�n, apertura, cierre y borrado de la tuber�a, y emisi�n y recepci�n de mensajes. Las implementaciones reales de una tuber�a pueden diferir, pero todas las implementaciones usan PBP para enlazar la tuber�a a un extremo.

Los mensajes de consulta de este protocolo se env�an por un extremo de la tuber�a para encontrar a un par en el extremo de la tuber�a enlazado al mismo anuncio de la tuber�a. El mensaje de consulta podr� solicitar informaci�n para obtener informaci�n acerca de los pares enlazados a la tuber�a. El mensaje de consulta tambi�n puede contener, opcionalmente, un ID que, si est� presente, indica que s�lo el par especificado debe responder a la consulta.

En cu�nto a los mensajes de respuesta, �stos se env�an por cada par asociado a la tuber�a al par solicitante, es decir, a aquel que realiz� la consulta.


\subsection{Protocolo de Encaminamiento al Destino (\textit{ERP - Endpoint Routing Protocol})}

El protocolo de encaminamiento al destino permite a los pares enviar mensajes a pares remotos sin tener una conexi�n directa con ellos. El mensaje se pasar� a trav�s de pares intermediarios para llegar a su destino final. ERP define un protocolo de consulta y respuesta que se utiliza para descubrir informaci�n de enrutamiento de pares, adjuntando a los mensajes utilizados en JXTA un mensaje especial que describe la ruta por la que debe viajar la informaci�n desde un par (el origen) a otro (el destino). 

Para enviar un mensaje a otro par, el origen primero busca en su memoria cach� para determinar si tiene una ruta hacia el par destino. Si no tiene una ruta hacia el destino, env�a a una consulta de solicitud de ruta pidiendo informaci�n de la ruta a seguir para el par de destino. Cualquier par que reciba esta consulta de ruta comprobar� si conoce alguna una ruta para el par solicitante. Si conoce una ruta v�lida, responder� a la consulta de ruta con la informaci�n de la ruta que �l posee para ese destino. De este modo, cualquier par puede realizar consultas sobre rutas y cualquier par dentro de un grupo de pares puede ofrecer informaci�n de ruta o distribuir mensajes destinados a otros pares. 

Las solicitudes de consulta de ruta son enviadas por un par origen para solicitar la informaci�n de ruta hacia un par de destino. Las respuestas de la ruta incluyen el ID del par de origen, el ID del par de destino de la ruta y una secuencia semi-ordenada de IDs de pares. Esta secuencia representa los pares que van a formar la ruta para transmitir el mensaje desde el par origen al par de destino, es decir, son los diferentes saltos o nodos que va a seguir el mensaje. Cada par a lo largo de la ruta indicada en la secuencia puede mejorar y optimizar la ruta bas�ndose en su propia informaci�n. Por ejemplo, si un par recibe un mensaje que contenga una ruta con diez saltos pero ese par est� directamente conectado al destino, puede reenviar el mensaje directamente, en lugar de enviarlo a lo largo de la ruta. 

El procedimiento de enrutamiento de mensajes utilizado por este protocolo sigue aproximadamente estos pasos:

\begin{enumerate}
	\item El par de origen comprueba si est� conectado al destino. Si no tiene una ruta hacia el destino, realiza una consulta de ruta y espera una respuesta, renunciando a la comunicaci�n si no la recibe en un tiempo determinado.
	\item El mensaje se difunde por los pares del grupo, estableciendo la secuencia de pares que pueden transmitir el mensaje hacia el destino pasando por otros pares. As�, se crea la ruta entre el origen y el destino.
	\item Si es un par intermedio en la ruta y no tiene una ruta hacia el destino o no se puede reenviar el mensaje a otro par que se haya establecido en la ruta encontrada, se manda un mensaje de fallo al origen, indicando que es imposible enviar su mensaje al destino.
	\item Si no ocurre la situaci�n anterior, el par intermedio lo reenv�a al siguiente par en la ruta, o al destino si est� conectado directamente con �ste.
	
\end{enumerate}


\subsection{Protocolo de Citas (\textit{RVP - Rendezvous Protocol})}

El protocolo de citas se utiliza para la propagaci�n de mensajes dentro de un grupo de pares. El RVP proporciona mecanismos para habilitar la propagaci�n de mensajes, realiz�ndose de un modo controlado y eficiente. Este protocolo se divide en tres partes:

\begin{milista}
	\item El protocolo utilizado por los pares de tipo \textit{Rendezvous} para organizarse, tambi�n conocido como protocolo \textit{\textbf{PeerView}}.
	\item El protocolo utilizado por pares cliente para registrar su inter�s de recibir mensajes de propagaci�n, conocido como protocolo simple de arrendamiento.
	\item El protocolo utilizado para la propagaci�n de mensajes a los pares que han registrado su inter�s para recibir mensajes. Este protocolo es el �nico que todos los participantes deben implementar.
 
 \end{milista}

\clearpage
\section{Iniciaci�n a la programaci�n de aplicaciones con JXTA}

El objetivo de esta secci�n es que sirva al lector como gu�a de referencia en su iniciaci�n a la programaci�n de aplicaciones basadas en JXTA. Para ello se comenzar� configurando el entorno de Java y se dar�n unas nociones b�sicas de configuraci�n del entorno de JXTA. Acto seguido, se indicar� como ejecutar el \textit{shell} de JXTA y, por �ltimo, se har� un peque�o ejemplo tipo \textit{HelloWorld}.

\subsection{Requisitos del sistema}

La implementaci�n actual de JXTA requiere que el sistema operativo tenga instalado una versi�n igual o posterior a la 5.0 del entorno de ejecuci�n de Java (JRE o \textit{Run-Time Environtment}) o del entorno de desarrollo (SDK o \textit{Software Development Kit}). Tanto el JDK como el JRE est�n disponibles para sistemas basados en Solaris, Microsoft Windows, Linux, Mac OS X, etc�tera.

\subsection{Instalaci�n y configuraci�n de los entornos}
\subsubsection{Java}
Para comprobar si ya dispone de una instalaci�n del entorno de ejecuci�n de Java en el equipo y que est� configurado correctamente, abra una ventana de comandos \textit{(Inicio -\textgreater Ejecutar -\textgreater cmd)} y escriba "`java -version"'. Debe obtener un resultado similar al de la figura \ref{fig:java-version}:

\imagen{./imagenes/java-version.png}{0.80}{C�mo consultar la versi�n de Java instalada}{fig:java-version}

Si aparece un mensaje de error, pruebe lo siguiente: 
\begin{milista}
	\item Compruebe si tiene instalado Java. Para ello compruebe que exista el directorio Java en \textit{C:\textbackslash Archivos de programa} (t�picamente). Si lo est�, config�relo como se indica m�s abajo.
	\item Si no existe, vaya a la siguiente direcci�n de Internet y descargue la �ltima versi�n del  Java SE Development Kit (JDK) \cite{download-jdk}
\end{milista}

Una vez descargado, ejecute el archivo, elija el directorio de instalaci�n y siga los pasos para completar la instalaci�n. Tras instalarlo, debe configurarlo. Para ello vaya a \textit{Inicio} y haga clic derecho en \textit{Mi PC}. En las opciones, elija la pesta�a de \textit{Opciones Avanzadas} y haga clic en \textit{Variables de entorno}. En la nueva ventana que aparece, busque la variable \textit{Path} y a�ada al final \texttt{;C:\textbackslash Archivos de programa\textbackslash Java\textbackslash jdk1.x.x\_xx\textbackslash bin}, donde "`x"' hace referencia a la versi�n que tiene instalada de Java. Ese es el directorio de instalaci�n por defecto de Java. Si en la instalaci�n eligi� otro diferente, debe introducir el que eligi� al instalarlo.

Una vez realizados los pasos descritos, pruebe a ejecutar de nuevo la sentencia "`java -versi�n"' en una ventana de comandos.


\subsubsection{JXTA}
Para comenzar, debe hacerse de una versi�n funcional de la librer�a de JXTA. Se proponen dos m�todos (ambos requieren que visite la p�gina de descargas de JXTA \cite{download-jxta}):
\begin{milista}
	\item Descargue el fichero \textit{jxse-src-2.5.zip}. Descompr�malo y comp�lelo con \textit{ant} o \textit{make}. Acto seguido se generar� el fichero \textit{jxta.jar} en el directorio \textit{./jxse-src-2.5/dist/}.
	\item Descargue el fichero \textit{jxse-lib-2.5.zip}. En su interior se encuentra el fichero \textit{jxta.jar} ya compilado.
\end{milista}

Tambi�n puede descargar la documentaci�n de JXTA (fichero \textit{jxse-doc-2.5.zip}) o la gu�a de programaci�n de JXTA para Java SE/EE (JXSE) 2.5 \cite{jxta-programming-guide}.

En versiones anteriores era necesario realizar una instalaci�n de dicho motor, pero hoy en d�a basta con incluir la librer�a de JXTA en la aplicaci�n y realizar una configuraci�n del \textit{peer}, bien sea en el c�digo fuente de la aplicaci�n o mediante un \textit{configurator} gr�fico externo (ver secci�n \ref{howto-run-jxta-apps}).

\subsection{C�mo ejecutar aplicaciones JXTA} \label{howto-run-jxta-apps}
Para ejecutar una aplicaci�n JXTA es necesario indicar donde se encuentran las librer�as, bien incluy�ndolas en el \textit{classpath} del sistema operativo o a�adi�ndolas al proyecto, en el entorno de programaci�n.

JXTA incluye una interfaz gr�fica (ver Figura \ref{fig:jxta-configurator}) que nos permite configurar el \textit{peer}, aunque normalmente este paso se suele evitar ya que cada aplicaci�n lo configura por c�digo utilizando la clase \textit{NetworkConigurator} o el \textit{NetkworkManager}, que abstrae la configuraci�n a una serie de configuraci�nes predefinidas:
\begin{milista}
	\item \textbf{Ad-Hoc}: un nodo t�picamente desplegado en una red ad-hoc, es decir, una red en la que no hay un nodo central y todos los dispositivos est�n en igualidad de condiciones. El \textit{par} no utilizar� una infraestructura de \textit{pares} (servicios \textit{rendezvous} o \textit{relay}).
	\item \textbf{Edge}: este tipo de nodos se pueden unir a una infraestructura de \textit{pares} (servicios \textit{rendezvous} o \textit{relay}).
	\item \textbf{Rendezvous}: proporiona servicios auto-suficientes (descubrimiento, etc.).
	\item \textbf{Relay}: proporciona servicios de mensajer�a, permitiendo atravesar corta-fuegos.
	\item \textbf{Proxy}: proporciona servicios de proxy.
	\item \textbf{Super}: proporciona la funcionalidad de un nodo \textit{rendezvous}, \textit{relay} o \textit{proxy}.
\end{milista}

\imagen{./imagenes/jxta-configurator.png}{0.3}{JXTA Configurator}{fig:jxta-configurator}

Una vez que el \textit{peer} ha sido configurado por alguno de los m�todos mencionados, se crear� autom�ticamente un directorio llamado \textit{.jxta} en el directorio de la aplicaci�n. Si se inspecciona dicho directorio, podr� encontrar un fichero llamado \textit{PlatformConfig}, que es el fichero de configuraci�n del par. Si desea volver a configurar el \textit{par}, se recomienda que borre dicho directorio.


\subsection{C�mo utilizar el \textit{Shell} de JXTA}
Para familiarizarse con JXTA, se recomienda practicar primero con el \textit{shell} que incorpora JXTA. De este modo se puede comprender el funcionamiento del sistema antes de adentrarse en el mundo de la programaci�n. A continuaci�n se explicar� un sencillo ejemplo de c�mo descubrir los pares del grupo por defecto, pero primero debe descargarse el shell de su p�gina Web \cite{download-shell}.

Una vez que haya descargado el shell, descompr�malo y ejecute el script \textit{run-shell.sh} (Linux/MacOS) o \textit{run-shell.bat} (Windows). Estos scripts incluyen autom�ticamente en su classpath la ubicaci�n de las librer�as de JXTA. Obtendr� una ventana como la que se muestra en la Figura \ref{fig:jxta-shell}. Si introduce la �rden \textit{'man'} obtendr� una lista con todas las �rdenes que puede ejecutar en el \textit{shell} junto a una breve descripci�n de cada una de ellas.

\imagen{./imagenes/jxta-shell.png}{0.3}{Shell de JXTA}{fig:jxta-shell}

\paragraph{Ejemplo pr�ctico: Descubrimiento de pares}
A continuaci�n se realizar� un sencillo ejemplo de c�mo descubrir los pares que hay en el grupo. Se propone al lector que inicie un \textit{shell} en dos ordenadores distintos, pertenecientes a una misma red, y ejecute las �rdenes "`peers"' y "`peers -r"'. Para obtener m�s informaci�n acerca de la orden "`peers"' puede ejecutar "`man peers"'. Si observa la figura \ref{fig:jxta-shell-peers} para ver los resultados de este ejemplo.

\imagen{./imagenes/jxta-shell-peers.png}{0.4}{Ejemplo de descubrimiento de pares}{fig:jxta-shell-peers}

Se propone al lector que haga uso de las �rdenes \textit{talk} y \textit{xfer} para establecer una conversaci�n de chat entre dos pares y para transferir ficheros, respectivamente. Utilice "`man talk"' y "`man xfer"' para obtener sus respectivas ayudas.

Por �ltimo, si no consigue que funcione alguno de los ejercicios propuestos, pruebe a deshabilitar el \textit{firewall} de su equipo.


\subsection{Ejemplo de un HelloWorld}
A continuaci�n se muestra el c�digo Java correspondiente a un programa de tipo \textit{Hola Mundo}. En este ejemplo se hace uso del \textit{NetworkManager} para configurar el \textit{par} en c�digo, por lo que en ning�n momento se mostrar� el \textit{JXTA UI Configurator}.

La sentencia \texttt{manager = new NetworkManager(ConfigMode.EDGE, "HelloWorld");} configura el \textit{par} con una configuraci�n por defecto: nodo de tipo \textit{edge} y le asigna el nombre \textit{HelloWorld}. Se podr�a personalizar esta configuraci�n obteniendo el objeto \textit{NetworkConfigurator} que encapsula el \textit{NetworkManager} con una sentencia del tipo a \texttt{NetworkConfigurator configurator = manager.getConfigurator();}.

El cometido de la aplicaci�n es sencillo:
\begin{enumerate}
	\item Crea una instancia del \textit{NetworkManager} e inicia la red JXTA.
	\item Se pone a la escucha de conexiones durante 2 minutos.
	\item Detiene la red JXTA.
\end{enumerate}

Una vez que la ejecuci�n de la aplicaci�n termina, se puede comprobar el directorio de configuraci�n que se ha creado (\textit{.jxta}):
\begin{milista}
	\item \textbf{PlatformConfig}: el fichero de configuraci�n creado por el \textit{NetworkManager}.
	\item \textbf{cm}: es un directorio que contiene un subdirectorio para cada grupo que se ha descubierto. En este ejemplo, deber�an aparecer los subdirectorios \textit{jxta-NetGroup} y \textit{jxta-WorldGroup}. Estos subdirectorios deben contener ficheros de �ndice (*.idx) y ficheros de anuncios (advertisements.tbl).
\end{milista}

\texttt{\lstinputlisting[inputencoding=latin1]{HelloWorld.java}}

\paragraph{Notas de inter�s}
\begin{enumerate}
	\item Cada vez que se instancia el NetworkManager, se crea un nuevo identificador de \textit{par} (JXTA PeerID), a menos que se haya cargado una configuraci�n local.
	\item Si se necesita cambiar la configuraci�n local de un \textit{par} (de EDGE a RENDEZVOUS, por ejemplo), se debe hacer mediante el NetworkManager utilizando el m�todo \textit{setMode()} \textbf{antes} de que se inicie la red.
\end{enumerate}

\clearpage
\section{Toolkits e implementaciones alternativas}
Aunque la implementaci�n de JXTA en Java es una buena herramienta para desarrollar aplicaciones P2P, no es la �nica alternativa de la que disponen los desarrolladores. Existen numerosas alternativas tanto de JXTA para otros lenguajes, como otra plataformas diferentes a JXTA.

\subsection{IBM BabbleNet}
\textit{AlphaWorks} es un grupo de IBM que desarrolla y distribuye nuevas tecnolog�as, principalmente para desarrolladores. BabbleNet \cite{babblenet} es una de estas tecnolog�as y consiste en un programa P2P descentralizado que permite a sus usuarios construir un chat en tiempo real sin la necesidad de conectarse a un servidor. El sistema est� desarrollado sobre un framekwork P2P escrito en Java que soporta comunicaciones entre nodos.

El c�digo fuente puede examinarse una vez instalado el software. La documentaci�n es m�nima y los comentarios que hay en el c�digo describen el sistema. Se trata de un software en fase experimental y con licencia \textit{open-source}.

\subsection{Intel}
Intel tambi�n se encuentra dentro del mercado del P2P, y lo hace con su \textit{Peer-to-Peer Acelerator Kit} \cite{p2p-acelerator-kit}. Se trata de un \textit{middleware} que se utiliza con Microsoft .NET.

\subsection{Microsoft .NET}
Microsoft ha desarrollado una p�gina Web en la que se muestra c�mo se pueden desarrollar sistemas P2P con la plataforma .NET \cite{net-p2p}. Algunas de las caracter�sticas que destacan son la siguientes:
\begin{milista}
	\item Todo el c�digo P2P se basa en la plataforma .NET.
	\item Los mensajes que se env�an entre los \textit{pares} son serializados como XML.
	\item Los \textit{pares} pueden compartir y acceder a distintos objetos.
	\item Tambi�n se ha implementado un servicio de descubrimiento utilizando .NET.
\end{milista}

\subsection{Peer-to-Peer Trusted Library}
\textit{Peer-to-Peer Trusted Library} (PtPTL) \cite{ptptl} es una biblioteca \textit{open-source} cuyo objetivo es profundizar en la seguridad de los sistemas P2P. Est� disponible tanto para sistemas basados en Windows como en sistemas basados en Linux. Sus principales caracter�sticas son las siguientes:
\begin{milista}
	\item Certificados digitales.
	\item Autenticaci�n de pares.
	\item Almacenamiento seguro.
	\item M�todos de codificaci�n sim�tricos y asim�tricos.
	\item Firmas digitales.
	\item Tarjetas digitales.
\end{milista}

Es importante apreciar que esta librer�a no es un \textit{toolkit} para desarrollar sistemas P2P, si no que es una biblioteca que permite dotar de "`confianza"' a dichos sistemas.


% Secci�n de ap�ndices
% Se a�ade el titulo de "Ap�ndices" al �ndice
%\addcontentsline{toc}{section}{Ap�ndices}

%\input{Apendices.tex}

% A�adimos la bibliografia al �ndice
\clearpage\phantomsection
\addcontentsline{toc}{section}{Referencias}

% Bibliograf�a. En este caso se usa BibTeX
\bibliographystyle{plain}
\pagestyle{plain} 
\bibliography{Bibliografia}

\end{document}
